\section{Impulse and Instinct: The Earth Race}

Prior to the lunar or Demetrian cultural stage and the patriarchal, or solar, cultural stage, \textbf{J J Bachofen} identified, through the study of ancient myths, and even earlier stage which he called ``hetaerism''. \textbf{Julius Evola} gives the name ``Telluric Race'' in the Sintesi to the guiding spirit of that age. (Full text follows below)

In such periods, marriage is unknown and women held in common. Bachofen explains:

\begin{quotex}
Woman is not endowed with all her charms in order to grow old in the arms of one man: the law of matter rejects all restrictions, abhors all fetters, and regards exclusivity as an offense against its divinity.
\end{quotex}


Unlike the Demetrians, who developed agriculture, the Tellurians were associated with the natural flora and fauna of the marshy lowlands, as well as swamp cults. This life is represented by witches, who were involved with herb collecting and concoctions using lower life forms like toads and lizards. The battle between the Telluric and Solar impulses is demonstrated, for example, in the Nordic legend of the Virgin Spring\footnote{\url{https://gornahoor.net/?p=1416}}, with Ingeri as the exemplar of the former and Karin of the latter.

The struggle between the Titans and the Olympic gods can also be understood in this context. The mlecchas were a degraded Telluric warrior caste, which the avatar Parashurama\footnote{\url{https://en.wikipedia.org/wiki/Parashurama}} came to destroy.

\paragraph{The Telluric Race}


The third race of the spirit that can be identified on the basis of ancient traditional symbols is the ``telluric'' or ``titanic'' race. It is a way of being that reveals itself by the adherence to life in all its immediacy, instinctiveness, and irrationality. The mind is today instinctively led to think of seismic phenomena by the term ``telluric'', more than recalling its etymology (from the Latin
\emph{tellus}, which means the earth), an assimilation that, in some

respect, could even have some justification. The telluric race is that of an explosive impulsivity, from sudden changes or from absolute identifications. As intensive as it is, it is equally gloomy, however without the depth and detachment necessary to also be tragic. Sexuality has a considerable part in it in its most elementary aspect: naturally, not just phallic, virile sexuality. In this regard, apart from truly inferior races, it can actually be said that it can rather more easily appear in a woman than in a man, in accordance with an entirely telluric nature. The feeling of the personality in the telluric man is poorly developed, the collective element predominates, in the sense that the ties of blood are always manifested in him in a material, atavistic, fatalistic form, something that can be clearly recognized in some typical aspects of the feeling of race and blood characteristic of the Hebraic people. In its appearance, not in primitivistic stages, but within a civilization previously formed by other human types, tellurism bears witness to the phase of the final decomposition of this civilization: it corresponds to the liberation and the unleashing, in a state of new freedom, of forces previously restrained by a higher law.

Through the appearance of sudden upheavals, one can recognize a telluric element in the race called ``desert'' by some race theorists, and also of a particular inner instability in the ``East Baltic'' race,. Through its gloomy and fatalistic side, the telluric man is then recognizable in the Etruscan race, in light of Bachofen's masterful description of it. Of course, the Mediterranean man generally, where he wants to form his own life according to a Nordic-Aryan style, has much to combat against this telluric possibility, even today. We then note that the telluric attribute was used by Keyserling, and not wrongly, to indicate an indisputable aspect of the so-called contemporary ``world revolution''.

In the cycles of primordial traditions, the titanic race appears as the natural antithesis to the race of Demetrian man, at the point where the original solar synthesis was lost. There we must then consider, above all, the degradation of the virile quality that now appears in a terrestrial form and makes its own not only the forms of a savage and violent assertion, but likewise some elemental forces of lower nature, tied, in times past, to the symbolism and the cult of Poseidon, for example. And one could in this regard, speak even of the promethean race, in as much as another distinctive trait of certain characteristics of such a race of the spirit is the attempt to usurp the dignity originally possessed by the solar race: from this arise the well-known myths of the battle of the Titans against the Olympic forces, or the memories contained in the Indo-Aryan tradition in regard to the mlecchas, the warrior race humiliated in the immense revolt by Parashurama, the representative of the oldest and highest spirituality when the ancestors of the Aryan conquerors of prehistoric India were still inhabiting the Hyperborean region.


\flrightit{Posted on 2015-06-14 by Aeneas}

\begin{center}* * *\end{center}

\begin{footnotesize}
\begin{sffamily}
\texttt{Mark Citadel on 2015-06-30 at 10:30 said:}

Would these societies have been highly primitive? I cant imagine any kind of actual civilization emerging from this.

\hfill

\texttt{Cologero on 2015-07-02 at 07:30 said:}

Yes, they are degenerate versions of Demetrian societies. What sort of ``civilization'' do you think will arise from movements like Femen or FreeTheNipple?

\end{sffamily}
\end{footnotesize}

